\documentclass[11pt]{article}

% Page layout: adjust margins slightly if needed to stay within 2 pages.
\usepackage[margin=0.7in]{geometry}
\usepackage{setspace}
\singlespacing
\usepackage{enumitem}
\usepackage{booktabs}
\usepackage{amsmath}
\usepackage[T1]{fontenc}
\usepackage{lmodern}
\usepackage[numbers,sort&compress]{natbib}

% Compact formatting to help meet the 2-page limit.
\setlength{\parindent}{0pt}
\setlength{\parskip}{3pt}
\setlist{nosep,leftmargin=*}
\usepackage{titlesec}
\titlespacing*{\section}{0pt}{5pt}{2pt}
\titleformat{\section}{\normalsize\bfseries}{\thesection.}{0.5em}{}

\title{Statistical Analysis of Causal Abstraction of Neural Networks}
\author{Dhruv Puri (MS Computer Science) \and
        Nichols Crawford Taylor (PhD Computer Science) \and
        Likhith Mallavarapu (MS Data Science)}
\date{CS 7200 Project Proposal, October 16, 2026}

\begin{document}
\maketitle
\vspace{-2em}

\section{Description of the Problem}
Distributed alignment search (DAS) tests whether a neural network implements a
high-level algorithm by learning a rotation of a hidden layer and measuring
interchange intervention accuracy (IIA): the share of input pairs on which the
intervened network and algorithm agree \citep{geiger2021causal, geiger2024finding}.
High IIA is read as evidence for the algorithm, but this test is uncalibrated: the main
results report IIA on the fitting data, as the best of several runs, without intervals,
and IIA exceeds chance even on random networks \citep{geiger2024finding}. Whether such
results can mislead is debated \citep{makelov2023subspacelookingforinterpretability,
sutter2025nonlinear}.

We ask how often the rule ``certify if IIA $\geq c$'' certifies an algorithm the
network does not implement, and whether held-out evaluation, null calibration and
selection correction fix this. The outcome is the certification decision (giving type
I error and power), the main predictor is the reporting procedure, and the system is
DAS applied to small networks trained on synthetic tasks with a known algorithm.

\section{Summary of the Data}
All data are simulated so the truth is known. We will use one or two tasks from the
causal abstraction literature that have a simple high-level algorithm and can be
generated synthetically; our leading candidate is hierarchical equality
\citep{geiger2024finding}, where four random vectors $w, x, y, z$ are labelled true
when $w = x$ and $y = z$ have the same truth value, and the final choice will follow a
pilot. For each task, networks of varying width are trained to perfect accuracy; known
positives are trained to compute the algorithm's variables \citep{geiger2022inducing},
and known negatives are random or follow a different algorithm. The unit of
observation is one DAS analysis of one network, yielding an IIA and a decision.
Predictors are the procedure, hidden width, subspace size relative to width, and
number of fitting pairs.

Input pairs sharing a base input are correlated, so binomial intervals are too narrow.
The procedures are applied to the same fits, so comparisons are paired. IIA clusters
near chance or 1, so variances are unequal. Non-converged fits will be excluded and
reported. Random networks may be easier negatives than real failures.

\section{Background}
DAS results are reported as point estimates, and critiques have relied on examples and
theory rather than error rates \citep{makelov2023subspacelookingforinterpretability,
sutter2025nonlinear}. Related statistical work compares probes against control tasks
\citep{hewitt2019control} and frames circuit claims as hypothesis tests
\citep{shi2024hypothesis}. We add simulated estimates of the type I error and power of
IIA-based certification under four reporting procedures.

\section{Methods}
For network $j$ with truth $T_j \in \{0,1\}$ (known negative or positive) and procedure
$p$, the response is the decision $Y_{jp} \in \{0,1\}$ (certified or not). The
procedures are: (1) \emph{naive}, the largest in-sample IIA over seeds, layers and
subspace sizes, compared with $c$; (2) \emph{held-out}, the rotation fit on one set of
base inputs and IIA measured on disjoint base inputs; (3) \emph{null-calibrated}, a
permutation-style $p$-value $\hat p = (1 + \#\{b: \mathrm{IIA}^{(0)}_b \geq
\mathrm{IIA}\})/(1+B)$, where $\mathrm{IIA}^{(0)}_b$ come from the same pipeline run on
calibration negatives and on pairs with shuffled counterfactual labels, certifying if
$\hat p \leq 0.05$; (4) \emph{selection-corrected}, the configuration chosen on a
selection split and tested once on a third split, or compared with the null
distribution of the maximum IIA over all configurations. The key quantities are the
type I error $\alpha_p(\mathbf{x}) = P(Y=1 \mid T=0, \mathbf{x})$ and power
$\pi_p(\mathbf{x}) = P(Y=1 \mid T=1, \mathbf{x})$, where $\mathbf{x}$ holds
$\log_2$ width, subspace size over width, and $\log$ number of fitting pairs. In each
design cell we estimate $\alpha_p$ and $\pi_p$ as proportions with Wilson intervals and
test $H_0\!: \alpha_p \leq 0.05$ with an exact binomial test. Across cells, the primary
model is a logistic regression, $\mathrm{logit}\,P(Y_{jp}=1) = \beta_{0p} +
\boldsymbol{\beta}_p^\top \mathbf{x}_j$, fit separately for negatives and positives.
Each IIA is reported with a 95\% percentile
interval from a cluster bootstrap that resamples base inputs, so that correlated pairs
stay together.

Because all four procedures are applied to the same networks, procedures are compared
with McNemar's test on paired decisions and, as an alternative, a mixed-effects logistic
regression with a random intercept per network. The calibration negatives in (3) are disjoint from
the negatives used to estimate $\alpha_p$, so the test is not evaluated on the data that
calibrated it. Diagnostics include binned residuals, likelihood-ratio tests for interactions, and
checks for separation, likely when error rates are near 0 or 1, in which case we use
Firth's penalized logistic regression. We apply Holm's correction across cells and
report $\alpha_p$ and $\pi_p$ over a range of thresholds $c$. The main difficulty is
cost: each DAS fit takes minutes, which limits networks per cell (30 per truth status
gives a standard error near 0.04 at $\alpha = 0.05$); the pooled regression borrows
strength across cells, and the pilot sets the design size. Wrong algorithms that share
variables with the true one are partly correct, so we record the agreement of their
counterfactual labels as a covariate.

\section{Preliminary Results}
None yet. We will first reproduce two results of \citet{geiger2024finding} (near-perfect
IIA on a trained network; IIA above chance on wide random networks), time one DAS fit,
and compare cluster and naive bootstrap intervals for IIA.

\section{Project Timeline}
\begin{center}
\small
\begin{tabular}{@{}p{0.22\linewidth}p{0.72\linewidth}@{}}
\toprule
\textbf{Week / Dates} & \textbf{Milestone} \\
\midrule
Oct. 16--23 & Set up DAS code; reproduce the published result. \\
Oct. 24--30 & Train known positives and negatives at 3 widths; time DAS fits; fix the design. \\
Oct. 31--Nov. 6 & Implement procedures, bootstrap and null; run naive and held-out over the design. \\
Nov. 7--13 & First type I error estimates; interim report and presentation (Nov. 10). \\
Nov. 14--20 & Run null-calibrated and selection-corrected procedures over the design. \\
Nov. 21--27 & Logistic and mixed-effects models, McNemar tests, diagnostics; power analysis. \\
Nov. 28--Dec. 4 & Robustness: wrong-algorithm negatives, sensitivity to $c$; draft report. \\
Dec. 5--11 & Peer review (Dec. 10); final report and presentation (Dec. 11). \\
\bottomrule
\end{tabular}
\end{center}

{\small
\setlength{\bibsep}{1pt}
\bibliographystyle{plainnat}
\bibliography{references}
}

\end{document}
